\documentclass[UTF8,a4paper,11pt,openany]{ctexbook}
\usepackage{../common/statlearnbook}
\SLSetBookMetadata{第 4 册}{统计学习方法讲义 v2.6.0}{无监督学习与矩阵分解}
\begin{document}
\setcounter{page}{564}
\setcounter{chapter}{25}
\setcounter{figure}{0}
\pagestyle{headings}
\chaptermark{聚类方法}
% v2.3.1 figure UID: FIG-P445-01
% Proposed post-v2.3.1 polish for FIG-P445-01: protect key-label size and stroke hierarchy.
\tikzset{slfig-FIG-P445-01/.style={font=\fontsize{9.2pt}{11.0pt}\selectfont,every path/.append style={line cap=round,line join=round}}}
\begin{figure}[H]
\centering
\begin{tikzpicture}[slfig-FIG-P445-01,x=1.35cm,y=1.35cm,
  cluster band/.style={rounded corners=2.2pt,line width=.44pt,draw opacity=.30,fill opacity=.075},
  cluster label/.style={font=\bfseries\fontsize{9.6pt}{11.3pt}\selectfont},
  merge axis/.style={-{Stealth[length=2.0mm,width=1.18mm]},line width=.72pt},
  merge tick/.style={line width=.50pt},
  dendro branch/.style={draw=SLInk,line width=.98pt},
  cut line/.style={draw=SLGray!88!black,line width=.82pt,densely dashed},
]
  % Cluster bands lie behind branches and stop below the cut height 1.40.
  \path[cluster band,draw=SLBlue,fill=SLBlue] (-.35,.12) rectangle (1.35,1.34);
  \path[cluster band,draw=SLGold,fill=SLGold] (1.65,.12) rectangle (2.35,1.34);
  \path[cluster band,draw=SLTeal,fill=SLTeal] (2.65,.12) rectangle (4.35,1.34);
  \node[cluster label,text=SLBlue] at (.5,.25) {$C_1$};
  \node[cluster label,text=SLGold!88!black] at (2,.25) {$C_2$};
  \node[cluster label,text=SLTeal] at (3.5,.25) {$C_3$};

  \draw[merge axis] (-.65,.05)--(-.65,3.35)
    node[above,rotate=90,anchor=south,font=\fontsize{9.4pt}{11.2pt}\selectfont] {合并高度};
  \foreach \y in {0,1,2,3}{\draw[merge tick] (-.70,\y)--(-.60,\y) node[left=2pt,font=\fontsize{8.6pt}{10.2pt}\selectfont]{\y};}
  \foreach \x/\lab in {0/$x_1$,1/$x_2$,2/$x_3$,3/$x_4$,4/$x_5$}
    \node[font=\fontsize{9.4pt}{11.2pt}\selectfont,anchor=north] at (\x,0) {\lab};

  % Explicit merge heights: .65, .95, 2.10, 3.00.
  \draw[dendro branch] (0,.08)--(0,.65)--(1,.65)--(1,.08);
  \draw[dendro branch] (3,.08)--(3,.95)--(4,.95)--(4,.08);
  \draw[dendro branch] (3.5,.95)--(3.5,2.10)--(2,2.10)--(2,.08);
  \draw[dendro branch] (.5,.65)--(.5,3.00)--(2.75,3.00)--(2.75,2.10);
  \draw[cut line] (-.45,1.40)--(4.45,1.40);
  \node[font=\fontsize{9.2pt}{11pt}\selectfont,text=SLGray!88!black,anchor=west] at (4.50,1.40) {切割高度 $h_c=1.4$};
\end{tikzpicture}
\caption{树的纵坐标表示合并高度；在虚线高度切割时，虚线以下的连通分支分别成为一个类。}\label{fig:V4-C02-dendrogram}
\end{figure}

\Cref{fig:V4-C02-dendrogram}对应层次聚类知识点：在所示高度切割得到$\{x_1,x_2\}$、$\{x_3\}$和$\{x_4,x_5\}$三个类；改变切割高度只合并或拆分嵌套分支，不改变已经记录的合并顺序。
\SLAdvancedSection{\kMeans{}算法与收敛}{sec:V4-C02-S05}
\textbf{学习算法。}\kMeans{}交替执行最近中心分配和类内均值更新，并在每次循环保留目标轨迹、空类修复记录与停止原因。
\end{document}
